\section{Maximal localization of the two positive physical sources}
\label{sec:v56-maximal-layers}
A bound for the mass of a boundary source does not bound its density
at every roof. It can, however, control all interval averages at once,
outside one explicitly estimated set of roofs. We prove that statement
for the original incidence and clearance sources. The exceptional set
is chosen before a roof resolution or a bounded source insertion is
specified. No physical history is removed from the density.

Throughout this section $I=[t,t+h]$ has fixed length $h>0$. Its role
is to localize a finite measure; the subintervals used below have no
positive lower bound on their lengths. For a nonnegative integrable
function $f$ on $I$ define
\begin{equation}\label{eq:v56-maximal-definition}
 \mathcal M_I f(u)=\sup_{J\ni u}\frac1{|J|}
                   \int_J\one_I(v)f(v)\,dv,
\end{equation}
where $J$ ranges over all bounded intervals of positive length in
$\R$. The supremum is unchanged if rational endpoints are used and
limits of such intervals are admitted. This defines a Borel function
on $\R$. In particular it bounds the average over every $J\subset I$
containing $u$, including one-sided intervals with endpoint $u$.
The letters $B$, $\varepsilon$ and $\lambda$ denote respectively a
reconstruction band, a physical protection width and a density level;
none denotes the interval length.

\subsection{A covering estimate and its source-weighted form}
\begin{lemma}[One envelope for every interval]
\label{lem:v56-covering}
For $f\ge0$ in $L^1(I)$ and $\lambda>0$, the set
$O_\lambda=\{\mathcal M_I f>\lambda\}$ is open and
\begin{equation}\label{eq:v56-maximal-weak-one}
 |O_\lambda|\le\frac3\lambda\int_I f.
\end{equation}
If, in addition, $|\{u\in I:f(u)>L\}|\le D L^{-p}$ for
$L>0$ and $p>1$, then
\begin{equation}\label{eq:v56-maximal-weak-p}
 |O_\lambda|\le C_pD\lambda^{-p}.
\end{equation}
Consequently, for $1<q<p$ and $\delta=\int_I f$,
\begin{equation}\label{eq:v56-maximal-interpolation}
 \int_{\R}(\mathcal M_I f)^q
       \le C_{p,q}D^{(q-1)/(p-1)}
                          \delta^{(p-q)/(p-1)}.
\end{equation}
The zero function is interpreted separately if $D\delta=0$.
\end{lemma}
\begin{proof}
Extend $f$ by zero. An interval with average strictly greater than
$\lambda$ can be enlarged slightly and retain this inequality, by
absolute continuity of the integral. Such enlarged open intervals
cover $O_\lambda$, proving openness. Cover a compact subset by
finitely many witnessing intervals. Select one of greatest length,
delete all intervals intersecting it, and repeat. The selected
intervals are disjoint. Every deleted interval is contained in the
concentric triple of its selected interval. The compact subset
therefore has measure at most three times the sum of the selected
lengths, at most $3\int f/\lambda$. Inner regularity proves
\eqref{eq:v56-maximal-weak-one}. This is the elementary interval
covering proof of the uncentered maximal inequality.

Write $f=f\one_{\{f\le\lambda/2\}}+
 f\one_{\{f>\lambda/2\}}$. The first summand has maximal average
at most $\lambda/2$. The layer-cake identity and the weak tail give
\[
 \int_{\{f>\lambda/2\}}f
 \le D(\lambda/2)^{1-p}
       +D\int_{\lambda/2}^{\infty}L^{-p}\,dL
 \le C_pD\lambda^{1-p}.
\]
Apply \eqref{eq:v56-maximal-weak-one} to the second summand with
threshold $\lambda/2$. This proves \eqref{eq:v56-maximal-weak-p}.
Finally integrate the minimum of $3\delta/\lambda$ and
$C_pD\lambda^{-p}$ in the distribution formula for the $q$th
power. Splitting at a constant multiple of
$(D/\delta)^{1/(p-1)}$ proves
\eqref{eq:v56-maximal-interpolation}. The integral near zero is
finite since $q>1$, even though the maximal function need not have
an integrable first power on all of $\R$.
\end{proof}

\begin{lemma}[Mass of a small roof set under the complete source]
\label{lem:v56-source-weight}
Suppose $P\ge0$ on $I$ satisfies
$|\{P>L\}|\le D L^{-p}$ for $L>0$, $p>1$. For every measurable
$E\subset I$,
\begin{equation}\label{eq:v56-source-weight}
 \int_E P\le\frac{p}{p-1}D^{1/p}|E|^{1-1/p}.
\end{equation}
If $P=G+f+e$, $f\ge0$, $\|G\|_\infty\le C_G$ and
$\|e\|_\infty\le\eta$, there is also the distinct estimate
\begin{equation}\label{eq:v56-source-weight-sharp}
 \int_E P\le(C_G+\eta)|E|+\int_I f.
\end{equation}
\end{lemma}
\begin{proof}
Integrate $\min\{|E|,DL^{-p}\}$ over $L>0$, splitting at
$(D/|E|)^{1/p}$. This proves the first estimate, with the cases of
zero measure immediate. The second follows by integrating the exact
positive-remainder representation. It uses a bounded controlled
part; it is not a consequence of the weak endpoint alone.
\end{proof}

\subsection{Finite-count localization of physical heights}
Use the unchanged densities of Section~\ref{sec:v48-physical-remainder}
and put
\[
 P=m^2p_{n,R}(k,m,\cdot),\qquad
 V_{\varepsilon,j}=m^2b^{\varepsilon,j,1}_{n,k,m,R},\quad
 V_\varepsilon=V_{\varepsilon,\mathrm{inc}}+
                            V_{\varepsilon,\mathrm{clr}}.
\]
Thus $0\le V_{\varepsilon,j}\le V_\varepsilon\le P$.
The normalization is $\nu_R^*=c^{-1}\nu|_{Y_R^*}$, used once.
The exact $n,k,m$ and the half-open occupation at $0,\ldots,m-1$
are unchanged; clearance of flight $j$ remains read at collision $j+1$.

\begin{theorem}[Finite-count maximal boundary control]
\label{thm:v56-finite-maximal}
Let $p_c=145/144$ and $0<\varepsilon<\varepsilon_0$. For every
$m\ge2$, every exact label, and every translated interval $I$,
set
\begin{equation}\label{eq:v56-physical-bad-set}
 O_{\varepsilon,\lambda}
       =\{u\in\R:\mathcal M_I V_\varepsilon(u)>\lambda\}.
\end{equation}
With constants uniform in $m,R,n,k,t,\varepsilon,\lambda$,
\begin{align}
 |O_{\varepsilon,\lambda}|
 &\le C(1+h)\min\{
       \varepsilon^{1/16}/\lambda,\lambda^{-p_c}\},
                                      \label{eq:v56-finite-bad-length}\\
 \int_{I\cap O_{\varepsilon,\lambda}}P
 &\le C(1+h)\min\{1,
                  (\varepsilon^{1/16}/\lambda)^{1/145}\}.
                                      \label{eq:v56-finite-bad-source}
\end{align}
On $I\setminus O_{\varepsilon,\lambda}$ both physical components
have height at most $\lambda$ almost everywhere, and
\begin{equation}\label{eq:v56-every-interval}
 \sup_{J\subset I:\,u\in J}\frac1{|J|}
                 \int_J V_{\varepsilon,j}\le\lambda
       \quad(u\in I\setminus O_{\varepsilon,\lambda}).
\end{equation}
The supremum is over all positive lengths. The same set works for
every bounded source insertion: the absolute density and every
interval average for $|w|\le M$ are bounded by $M$ times these
quantities. Moreover, for $1<q<p_c$,
\begin{equation}\label{eq:v56-maximal-physical-Lq}
 \int_{\R}(\mathcal M_I V_\varepsilon)^q
       \le C_q(1+h)\varepsilon^{9(p_c-q)}.
\end{equation}
\end{theorem}
\begin{proof}
Theorem~\ref{thm:v55-error-free-mass} gives
$\int_I V_\varepsilon\le C(1+h)\varepsilon^{1/16}$,
simultaneously for every finite count and protection scale.
Corollary~\ref{cor:v55-physical-endpoint} gives the weak $L^{p_c}$
bound for $P$ and hence for $V_\varepsilon$. Apply
Lemma~\ref{lem:v56-covering} to get
\eqref{eq:v56-finite-bad-length}. Lemma~\ref{lem:v56-source-weight}
then gives \eqref{eq:v56-finite-bad-source}, since
$1-1/p_c=1/145$. The uniform first moment of $P$ supplies the
alternative bound by $C(1+h)$.

Lebesgue differentiation gives
$V_\varepsilon\le\mathcal M_I V_\varepsilon$ almost everywhere
on $I$. This proves the height assertion. The interval assertion
is the definition of the maximal function and needs no exceptional
null set depending on an interval. Positive source domination gives
the inserted assertions. More precisely, fix a regular disintegration
of the finite physical source over its roof. Defining every inserted
density by this one kernel gives
$|V_{\varepsilon,j}^w(u)|\le M V_{\varepsilon,j}(u)$ in that
common version. A separately chosen density version for each $w$
is not being used. The last bound follows from
\eqref{eq:v56-maximal-interpolation}, with
$\delta=C(1+h)\varepsilon^{1/16}$, $D=C(1+h)$ and
$(1/16)/(p_c-1)=9$. All these are finite inequalities; no spectral
band is made to depend on a roof resolution.
\end{proof}

\subsection{A sharper source estimate in the ordered raw limit}
For large fixed $B$, take $\varepsilon(B)=A_0B^{-1/12}$ and write
$V_B=\mathfrak b_{B,m,R}^{n,k}$. Define
\[
 \eta_{B,m}=\sup_{R,n,k}\|P-G-V_B\|_\infty,
 \qquad G(u)=\mathcal L_{m,R}(k_1,k_2,u,n).
\]
It is finite for each count, and
$\limsup_m\eta_{B,m}\le CB^{-1/192}$ by
Theorem~\ref{thm:v51-positive-error}. The parameter-uniform
arithmetic kernel is bounded by a fixed $C_G$. Let
\begin{equation}\label{eq:v56-band-threshold}
 \lambda_B=B^{-1/384},\qquad
 O_{B,m,R}^{n,k,I}=O_{\varepsilon(B),\lambda_B}.
\end{equation}

\begin{corollary}[Resolution-uniform localization of the raw error]
\label{cor:v56-ordered-localization}
For every fixed $h>0$,
\begin{align}
 \sup_{m,R,n,k,t}|O_{B,m,R}^{n,k,I}|
       &\le C(1+h)B^{-1/384},
                                      \label{eq:v56-ordered-length}\\
 \limsup_{m\to\infty}\sup_{R,n,k,t}
       \int_{I\cap O_{B,m,R}^{n,k,I}}P
       &\le C(1+h)B^{-1/384}.          \label{eq:v56-ordered-source}
\end{align}
For $u\in I\setminus O_{B,m,R}^{n,k,I}$,
\begin{equation}\label{eq:v56-maximal-raw-error}
 \sup_{J\subset I:\,u\in J}\frac1{|J|}\int_J|P-G|
                         \le\lambda_B+\eta_{B,m}.
\end{equation}
The pointwise bound $|P(u)-G(u)|\le\lambda_B+\eta_{B,m}$ also
holds almost everywhere on this same complement. Consequently all
these errors have ordered limsup at most $CB^{-1/384}$.
\end{corollary}
\begin{proof}
Equation~\eqref{eq:v56-finite-bad-length} gives
\eqref{eq:v56-ordered-length}. Apply
\eqref{eq:v56-source-weight-sharp} with $f=V_B$ to obtain the
finite estimate
\[
 \int_{I\cap O}P\le
 C(1+h)\{(C_G+\eta_{B,m})B^{-1/384}+B^{-1/192}\}.
\]
Its collision limsup proves \eqref{eq:v56-ordered-source}.
This is sharper than the finite-count weak-endpoint estimate
\eqref{eq:v56-finite-bad-source}, because it also uses the
positive-remainder raw representation. Finally
$|P-G|\le V_B+\eta_{B,m}$. Take interval averages at a good
anchor $u$ and use \eqref{eq:v56-every-interval}, then use
Lebesgue differentiation for the pointwise assertion.
\end{proof}

\paragraph{The part of the height problem still requiring geometry.}
The open set in \eqref{eq:v56-physical-bad-set} localizes all upward
physical concentration, not just concentration at one prescribed
resolution. Its length and original source mass tend to zero in the
stated ordered limit. The set need not be empty. A thin positive spike
can lie entirely inside it. Thus \eqref{eq:v56-maximal-raw-error}
retains the original arithmetic pointwise target but does not prove
\eqref{eq:v51-positive-criterion} on the excluded roofs. No incidence
or competing-hit seam has been traversed in this argument.
